\section{TreeSHAP và cấu trúc của mô hình cây}
\label{sec:ch04-treeshap}

\index{TreeSHAP}%
KernelSHAP vẫn phải gọi mô hình cho nhiều liên minh, nên nó phù hợp khi chỉ có
quyền truy cập \tn{hộp đen}{black box}. Với mô hình cây, cấu trúc các nút tách
cho phép tính đóng góp bằng quy hoạch động trên các đường đi của cây, thay vì
liệt kê từng liên minh. Tên TreeSHAP gắn với dòng thuật toán ấy trong thực
hành; bài báo gốc của SHAP không đặt tên này, mà đặt tên cho Kernel SHAP,
Linear SHAP, Low-Order SHAP, Max SHAP và Deep SHAP~\autocite{p10shap}. Các công trình sau
SHAP gốc mô tả các thuật toán quy hoạch động hiệu quả cho giá trị SHAP của mô
hình cây~\autocite{p19shapinference}.

Lợi ích ở đây là tính toán, không phải là một định nghĩa mới về đóng góp. Nếu
hai phương pháp dùng cùng hàm giá trị $v$, chúng phải nhắm tới cùng
\tn{giá trị Shapley}{Shapley value}. Điều làm phép tính rẻ đi là cấu trúc cây: một đường đi chỉ kiểm tra một số đặc
trưng, và nhiều liên minh có chung tiền tố.

Tuy nhiên, tốc độ không giải quyết lựa chọn về đặc trưng còn thiếu. Cây phải
đi theo nhánh nào khi một đặc trưng còn thiếu? Câu trả lời vẫn cần một quy tắc
hoàn thiện. Vì vậy, khi đọc một lời giải thích từ mô hình cây, cần tách
hai câu hỏi: thuật toán đã tính các liên minh hiệu quả đến đâu, và các liên
minh đó đã được gán nghĩa theo phân phối nào.
